% Part: first-order-logic
% Chapter: tableaux
% Section: derivations

\documentclass[../../../include/open-logic-section]{subfiles}

\begin{document}

\iftag{FOL}
      {\olfileid{fol}{tab}{der}}
      {\olfileid{pl}{tab}{der}}

\olsection{\usetoken{P}{tableau}}

\begin{explain}
పరికల్పన అంటే ఏమిటో చెప్పాం; నిగమన నియమాలనూ ఇచ్చాం. \tetoken{టాబ్లోలు}{!!^{tableau}s}
వీటి నుంచి ఆగమనాత్మకంగా ఉత్పన్నమవుతాయి: ప్రతి \tetoken{టాబ్లో}{!!{tableau}}, ఒకటి
లేదా అంతకంటే ఎక్కువ పరికల్పనలతో కూడిన ఒకే శాఖ అయినా అవుతుంది, లేదా
ఒక శాఖపై నిగమన నియమాల్లో ఒకదాన్ని వర్తింపజేసి మరొక \tetoken{ఒక టాబ్లో}{!!a{tableau}}
నుంచి ఏర్పడుతుంది.
\end{explain}

\begin{defn}[\tetoken{టాబ్లో}{!!^{tableau}}]
పరికల్పనలు $\sFmla{S_1}{!A_1}$, \dots,
$\sFmla{S_n}{!A_n}$ (ప్రతి $S_i$, $\True$ లేదా~$\False$లో ఒకటి)
కోసం ఒక \tetoken{టాబ్లో}{!!^a{tableau}} అంటే కింది షరతులను తీర్చే \tetoken{చిహ్నిత సూత్రాలు}{!!{signed formula}s}
పరిమిత వృక్షం:
\begin{enumerate}
\item వృక్షంలో అన్నిటికంటే పైనున్న $n$ \tetoken{చిహ్నిత సూత్రాలు}{!!{signed formula}s},
  ఒకదాని కింద మరొకటిగా $\sFmla{S_i}{!A_i}$లు.
\item పరికల్పనల్లో ఒకటి కాని వృక్షంలోని ప్రతి \tetoken{చిహ్నిత సూత్రం}{!!{signed formula}},
  దాని పైనున్న శాఖలోని ఒక \tetoken{చిహ్నిత సూత్రానికి}{!!a{signed formula}} నిగమన నియమాన్ని
  సరిగ్గా వర్తింపజేయడం వల్ల వస్తుంది.
\end{enumerate}
ఒక \tetoken{టాబ్లో}{!!a{tableau}} శాఖలో $\sFmla{\True}{!A}$,
$\sFmla{\False}{!A}$ రెండూ ఉంటే, అప్పుడు మాత్రమే అది
\emph{సంవృతమైనది}; లేకపోతే అది \emph{వివృతమైనది}. ప్రతి శాఖా
సంవృతంగా ఉన్న ఒక \tetoken{టాబ్లో}{!!^a{tableau}}, తన పరికల్పనల సమితికి ఒక
\emph{సంవృత \tetoken{టాబ్లో}{!!{tableau}}}. ఒక \tetoken{టాబ్లో}{!!a{tableau}} సంవృతం కాకపోతే, అంటే
అందులో కనీసం ఒక వివృత శాఖ ఉంటే, అది \emph{వివృతమైనది}.
\end{defn}

\begin{ex}
  ప్రతి పరికల్పనల సమితి తనంతట తానే ఒక \tetoken{టాబ్లో}{!!a{tableau}}; కానీ సాధారణంగా
  అది సంవృతంగా ఉండదు. (పరికల్పనల్లో ఇప్పటికే
  $\sFmla{\True}{!A}$,~$\sFmla{\False}{!A}$ అనే \tetoken{చిహ్నిత సూత్రాలు}{!!{signed formula}s}
  జత ఉంటే మాత్రమే అది సంవృతమని స్పష్టం.)

  ఒక \tetoken{టాబ్లోలోని}{!!a{tableau}} ఒక \tetoken{చిహ్నిత సూత్రం}{!!a{signed formula}}~$!A$కు నిగమన
  నియమాల్లో ఒకదాన్ని వర్తింపజేసి, ఆ టాబ్లో వివృతమైనదైనా
  సంవృతమైనదైనా, దాని నుంచి కొత్తదైన మరింత పెద్ద టాబ్లోను
  పొందవచ్చు. మనం నియమాన్ని వర్తింపజేసే $!A$ సంభవాన్ని కలిగి ఉన్న
  ప్రతి శాఖ చివరికి ఆ నియమం ఒకటి లేదా అంతకంటే ఎక్కువ
  \tetoken{చిహ్నిత సూత్రాలను}{!!{signed formula}s} చేర్చుతుంది.

  ఉదాహరణకు, $\sFmla{\True}{!A \land \lnot !A}$ అనే పరికల్పనను
  పరిశీలించండి. ఆ పరికల్పన మాత్రమే గల (వివృత) \tetoken{టాబ్లో}{!!{tableau}} ఇది:
  \begin{center}
    \begin{tableau}{}
      [\sFmla{\True}{\formula{A} \land \lnot \formula{A}}, just=\TAss]
    \end{tableau}{}
  \end{center}
  పరికల్పనకు $\TRule{\True}{\land}$ నియమాన్ని వర్తింపజేసి దాని నుంచి
  కొత్త \tetoken{టాబ్లోను}{!!{tableau}} పొందుతాం. ఆ నియమం \tetoken{టాబ్లోకు}{!!{tableau}}
  $\sFmla{\True}{!A}$, $\sFmla{\True}{\lnot !A}$ అనే రెండు కొత్త
  వరుసలను చేర్చనిస్తుంది:
  \begin{center}
    \begin{tableau}{}
      [\sFmla{\True}{\formula{A} \land \lnot \formula{A}}, just=\TAss
        [\sFmla{\True}{\formula{A}}, just={\TRule{\True}{\land}[1]},
          [\sFmla{\True}{\lnot \formula{A}}, just={\TRule{\True}{\land}[1]}
          ]
        ]
      ]
    \end{tableau}{}
  \end{center}
  \tetoken{టాబ్లోలను}{!!{tableau}s} రాసేటప్పుడు, మనం వర్తింపజేసిన నియమాలను కుడివైపు
  నమోదు చేస్తాం (ఉదాహరణకు, $\TRule{\True}{\land} 1$ అంటే ఆ వరుసలోని
  \tetoken{చిహ్నిత సూత్రం}{!!{signed formula}}, వరుస~$1$లోని \tetoken{చిహ్నిత సూత్రానికి}{!!{signed formula}}
  $\TRule{\True}{\land}$ నియమాన్ని వర్తింపజేయడం వల్ల వచ్చిందని
  అర్థం). ఈ కొత్త \tetoken{టాబ్లోలో}{!!{tableau}} ఇప్పుడు అదనపు \tetoken{చిహ్నిత సూత్రాలు}{!!{signed formula}s}
  ఉన్నాయి; కానీ వాటిలో ఒక్కదానికే ($\sFmla{\True}{\lnot !A}$కు)
  ఒక నియమాన్ని వర్తింపజేయగలం (ఈ సందర్భంలో
  $\TRule{\True}{\lnot}$ నియమం). దాని వల్ల కింది సంవృత
  \tetoken{టాబ్లో}{!!{tableau}} వస్తుంది:
  \begin{center}
    \begin{tableau}{}
      [\sFmla{\True}{\formula{A} \land \lnot \formula{A}}, just=\TAss
        [\sFmla{\True}{\formula{A}}, just={\TRule{\True}{\land}[1]}
          [\sFmla{\True}{\lnot \formula{A}}, just={\TRule{\True}{\land}[1]}
            [\sFmla{\False}{\formula{A}}, just={\TRule{\True}{\lnot}[3]}, close]
          ]
        ]
      ]
    \end{tableau}{}
  \end{center}
\end{ex}

\end{document}
